\section{Parallel repetition}
\label{ch:parallel-repetition}

This chapter states the parallel repetition theorem used for gap amplification in the
compression procedure (Section~\ref{sec:ps-parallel-repetition}):
\emph{direct} (unanchored) parallel repetition in value form, for entangled and for
commuting-operator strategies, together with Lin's tracial density theorem. Both
theorems are imported from external Lean developments vendored into this repository,
and the chapter records their statements in the vocabulary of
Section~\ref{ch:foundations} together with the bridges connecting the two; both are
formalized end to end, the entangled one through Lemma~\ref{lem:povm-value-eq}.

The \emph{anchored}, entanglement-preserving theorem of Bavarian, Vidick and
Yuen~\cite{BVY17}, which~\cite{JNVWY20} uses for this step, is not stated here.
Anchoring is what the entanglement form forces: \cite{MNY}'s recursive compression
lemma needs a measure $f$ with $f(\compr(\hat x, n)) \geq \max\{f(x), n\}$, so every
transformation has to preserve $\Ent$, and only the anchored theorem does. The
compressibility criterion of~\cite{Lin25} used here (Lemma~\ref{lem:compressible-criterion})
carries no measure at all, so value decay suffices and the repeated game may be the
game as it is. Remark~\ref{rem:direct-vs-anchored} records the comparison, and
Remark~\ref{rem:entanglement-form} what the entanglement clauses of~\cite{JNVWY20}
would buy.

\subsection{Direct parallel repetition}
\label{sec:pr-direct}

The $n$-fold direct repetition of a game asks $n$ independent question pairs at once and
accepts if all $n$ answer pairs are accepted. Two recent Lean developments prove that
the value of the repeated game decays exponentially in $n$, for every game, with no
anchoring and at a rate depending only on the gap $1 - \val(\game)$ and on the size of
the answer alphabets: one for tensor-product (entangled) strategies, by
OpenAI~\cite{OpenAI26}, and one for commuting-operator strategies~\cite{Vid26}, the
latter building on the former and on Lin's tracial density theorem~\cite{Lin23}. Both
are vendored into this repository (\texttt{MIPRE/Background/Repetition/}, with
provenance recorded in the \texttt{README.md} of each directory) and connected to the
definitions of Section~\ref{ch:foundations} by explicit bridges; the statements below
are the bridged forms. The commuting-operator theorem is recorded now, although only
the entangled one enters $\MIPstar = \RE$, because it is the form needed for the
commuting-operator counterpart $\MIPco = \coRE$~\cite{Lin25} (Remark~\ref{rem:further}).

\begin{definition}[Direct repetition]
\label{def:direct-repetition}
\uses{def:game,def:sync-game}
\lean{MIPRE.Game.repeat, MIPRE.SynchronousGame.repeat}
\leanok
Let $\game = (\mX, \mY, \mA, \mB, \mu, D)$ be a game and $n \geq 0$ an integer. The
\emph{$n$-fold direct repetition} of $\game$ is the game $\game^{n}$ with question
alphabets $\mX^n$, $\mY^n$ and answer alphabets $\mA^n$, $\mB^n$, question distribution
$\mu^{n}(\bar{x}, \bar{y}) = \prod_{i=1}^{n} \mu(x_i, y_i)$, and decision predicate
$D^{n}(\bar{x}, \bar{y}, \bar{a}, \bar{b}) = \bigwedge_{i=1}^{n} D(x_i, y_i, a_i, b_i)$.
The direct repetition of a synchronous game is synchronous, with the same data.
\end{definition}

The entangled theorem is stated for the quantum value of Definition~\ref{def:q-value},
whose strategies are bipartite. The commuting-operator theorem needs the bipartite
counterpart of the commuting value, which Section~\ref{ch:foundations} defines in
synchronous (tracial) form only (Definition~\ref{def:commuting-value}); the two are
related by Lemma~\ref{lem:tracial-le-co}.

\begin{definition}[Commuting-operator strategy]
\label{def:co-strategy-bipartite}
\lean{MIPRE.CommutingOperatorStrategy}
\leanok
\qlib{}
A \emph{commuting-operator strategy} for $\game$ consists of a Hilbert space $\mH$ (a
complete inner product space over $\C$, of any dimension), a unit vector $\ket{\psi}
\in \mH$, and POVMs $\{E^x_a\}_{a \in \mA}$ and $\{F^y_b\}_{b \in \mB}$ on $\mH$, one
for each $x$ (resp.\ $y$) --- each $E^x_a$ and $F^y_b$ is a positive bounded operator,
$\sum_a E^x_a = \Id$ and $\sum_b F^y_b = \Id$ --- such that $E^x_a F^y_b = F^y_b
E^x_a$ for all $x, y, a, b$. The players answer $(a, b)$ to questions $(x, y)$ with
probability $\bra{\psi} E^x_a F^y_b \ket{\psi}$.
\end{definition}

\begin{definition}[Commuting-operator value]
\label{def:co-value-bipartite}
\uses{def:co-strategy-bipartite,def:game}
\lean{MIPRE.CommutingOperatorStrategy.value, MIPRE.commutingOperatorValue}
\leanok
\qlib{}
The value of a commuting-operator strategy $\strategy = (\mH, \psi, E, F)$ in $\game$
is
\begin{equation*}
  \omegaco(\game, \strategy) = \sum_{x, y, a, b} \mu(x,y)\, D(x,y,a,b)\,
  \bra{\psi} E^x_a F^y_b \ket{\psi}\;,
\end{equation*}
and the \emph{commuting-operator value} $\omegaco(\game)$ of $\game$ is the supremum of
$\omegaco(\game, \strategy)$ over all commuting-operator strategies.
\end{definition}

\paragraph{Comments.} As in Definition~\ref{def:q-value}, the formalization takes the
real part of the (real, nonnegative) probability. The Hilbert space of a strategy lives
in the lowest universe, as in the vendored development; this is no restriction for the
suprema, which are bounded by $1$. The notation $\omegaco$ follows~\cite{Vid26} and
keeps the bipartite value apart from the tracial value $\valco$ of
Definition~\ref{def:commuting-value}.

\effortMedium\ \qlib

\begin{lemma}[Tracial strategies are commuting-operator strategies]
\label{lem:tracial-le-co}
\uses{def:sync-game,def:commuting-strategy,def:commuting-value,def:co-value-bipartite}
\qlib{}
For every synchronous game $\game$, $\valco(\game) \leq \omegaco(\game)$.
The Lean proof of the main theorem does not use this lemma; it serves the downstream
$\MIPco = \coRE$ chapter only.
\end{lemma}
\begin{proof}
Let $(\alg, \tau, M)$ be a commuting strategy (Definition~\ref{def:commuting-strategy})
and let $\mH = L^2(\alg, \tau)$ be the completion of $\alg$ for the inner product
$\langle a, b \rangle = \tau(a^* b)$, with the class of $1$ as the unit vector $\xi$.
Left multiplication $\pi(a) : b \mapsto ab$ and right multiplication $\rho(a) : b
\mapsto ba$ extend to bounded operators on $\mH$ (for $\rho$ this uses traciality:
$\norm{ba}^2 = \tau(a^* b^* b a) = \tau(b^* b\, a a^*) \leq \norm{a}^2 \norm{b}^2$),
commute with each other, and satisfy $\rho(p)^* = \rho(p^*)$ (again by traciality) so
that $\rho$ maps projections to projections. Setting $E^x_a = \pi(M^x_a)$ and $F^y_b =
\rho(M^y_b)$ gives a commuting-operator strategy with $\bra{\xi} E^x_a F^y_b \ket{\xi}
= \langle 1, M^x_a M^y_b \rangle = \tau(M^x_a M^y_b)$, hence with the same value as
$(\alg, \tau, M)$. Taking suprema yields the claim.
\end{proof}

\paragraph{Comments.} This is the direction needed to apply
Theorem~\ref{thm:direct-repetition-co} to the tracial value used in the pipeline: a
bound on $\omegaco(\game^n)$ is a bound on $\valco(\game^n)$. The reverse comparison
(from $\omegaco$ to $\valco$, false in general and true up to error for synchronous games
that carry proportional weight on the diagonal, as in
Remark~\ref{rem:almost-sync-hypothesis}) is the commuting-operator analogue of
Theorem~\ref{thm:almost-sync}, proved by Lin~\cite{Lin23}; it is part of the transport
discussed in Section~\ref{sec:departures} and not of this section. The Lean statement
of this lemma and its GNS construction are a sub-project of their own.

\effortHard\ \qlib

\begin{theorem}[Direct parallel repetition for commuting-operator strategies]
\label{thm:direct-repetition-co}
\uses{def:direct-repetition,def:co-value-bipartite}
\lean{MIPRE.Repetition.commutingOperatorValue_repeat_le}
\leanok
\qlib{}
There is a universal constant $c > 0$ such that for every game $\game$ with nonempty
alphabets and every integer $n \geq 1$, writing $\eps = 1 - \omegaco(\game)$,
\begin{equation*}
  \omegaco\bigl( \game^{n} \bigr) \;\leq\;
  \exp\Bigl( - \frac{c\, \eps^{7}\, n}{\eps + \log\bigl(\abs{\mA}\,\abs{\mB}\bigr)} \Bigr)\;.
\end{equation*}
\end{theorem}
\begin{proof}
\leanok
This is the main theorem of~\cite{Vid26} (Theorem~7.1 there), whose Lean proof is
vendored at \texttt{MIPRE/Background/Repetition/CommutingRepetition/}. The bridge
identifies a game of Definition~\ref{def:game} with the game of the vendored
development having the $\{0, 1\}$-valued payoff $D$, a commuting-operator strategy of
Definition~\ref{def:co-strategy-bipartite} with a strategy of the vendored development
(field by field, so that the two suprema defining the value range over the same set of
numbers), and the two direct repetitions with each other. When $\eps = 0$ and $\abs{\mA}
= \abs{\mB} = 1$ the formalization reads the quotient as $0$ and the bound is trivial.
\end{proof}

\begin{theorem}[Direct parallel repetition for entangled strategies]
\label{thm:direct-repetition-q}
\uses{def:direct-repetition,def:q-value,lem:povm-value-eq}
\lean{MIPRE.Repetition.quantumValue_repeat_le}
\leanok
\qlib{}
There is a universal constant $c > 0$ such that for every game $\game$ with nonempty
answer alphabets and $\valstar(\game) < 1$, and every integer $n \geq 1$, writing
$\eps = 1 - \valstar(\game)$,
\begin{equation*}
  \valstar\bigl( \game^{n} \bigr) \;\leq\;
  \exp\Bigl( - \frac{c\, \eps^{13}\, n}{\eps + \log\bigl(\abs{\mA}\,\abs{\mB}\bigr)} \Bigr)\;.
\end{equation*}
\end{theorem}
\begin{proof}
\leanok
This is the main theorem of Chapter~6 of~\cite{OpenAI26}, whose Lean proof is vendored
at \texttt{MIPRE/Background/Repetition/TenProofs/}. Its entangled value is defined with
a density matrix on the product of two finite sets and POVMs; Lemma~\ref{lem:povm-value-eq}
identifies it with $\valstar$ (Definition~\ref{def:q-value}), after which the two direct
repetitions are literally the same construction.
\end{proof}

\effortMedium\ \qlib

\begin{lemma}[Mixed states and POVMs do not change the quantum value]
\label{lem:povm-value-eq}
\uses{def:tensor-strategy,def:q-value}
\lean{MIPRE.Repetition.quantumValue_eq_entangledValue}
\leanok
\qlib{}
Let $\valstar_{\mathrm{POVM}}(\game)$ be the supremum of the winning probabilities
$\sum_{x,y,a,b} \mu(x,y)\, D(x,y,a,b)\, \tr\bigl( (A^x_a \ot B^y_b)\, \rho \bigr)$ over
all finite sets $S_\alice$, $S_\bob$, density matrices $\rho$ on $\C^{S_\alice} \ot
\C^{S_\bob}$ and POVMs $\{A^x_a\}$ on $\C^{S_\alice}$, $\{B^y_b\}$ on $\C^{S_\bob}$.
Then $\valstar_{\mathrm{POVM}}(\game) = \valstar(\game)$.
\end{lemma}
\begin{proof}
\uses{lem:naimark-dilation}
\leanok
A tensor-product strategy (Definition~\ref{def:tensor-strategy}) is a strategy of the
second kind, with $S_\alice = \{1, \ldots, d_\alice\}$, $S_\bob = \{1, \ldots,
d_\bob\}$, $\rho = \ket{\psi}\bra{\psi}$ and projective measurements, which gives
$\valstar \leq \valstar_{\mathrm{POVM}}$.

Conversely, given a strategy of the second kind, purify $\rho$ on a reference system
placed on Bob's side, where Bob's effects act as $B^y_b \ot \Id$, so that the
bipartite structure survives; then dilate the two POVM \emph{families} to projective
measurements on one extra register each, and reindex the two finite sets by initial
segments of $\N$. Neither step changes an outcome probability, so the result is a
tensor-product strategy of the same value.

The dilation is the only delicate point, and the reason is that the shared state may
not depend on the questions while the dilating isometry does. Writing $A^x_a =
(K^x_a)^\dagger K^x_a$, the isometry $V_x : v \mapsto \sum_a (K^x_a v) \ot \ket{a}$
satisfies $V_x^\dagger (\Id \ot \ket{a}\bra{a}) V_x = A^x_a$; extending each $V_x$ to a
unitary $U_x$ of $\C^{S_\alice} \ot \C^{\mA}$ and setting $P^x_a := U_x^\dagger (\Id
\ot \ket{a}\bra{a}) U_x$ gives projective measurements all of which are compressed back
to the POVMs by the \emph{same}, question-independent isometry $v \mapsto v \ot
\ket{a_0}$, so that the single state $\ket{\psi} \ot \ket{a_0} \ot \ket{b_0}$
works for every question. In Lean the extension is Mathlib's completion of an orthonormal
family to an orthonormal basis, and only the factorization $A^x_a = (K^x_a)^\dagger K^x_a$ is
used, not a matrix square root, which Mathlib does not have. The nonemptiness of $\mA$
and $\mB$ needed to choose $a_0$ and $b_0$ is derived, not assumed: $\mu$ is a
probability distribution, so some question is asked, and a POVM on a space carrying a
state has a nonempty outcome set.
\end{proof}

\paragraph{Comments.} The lemma is a general fact of independent interest, and the
substantive part of the entangled bridge; Naimark dilation is not in Mathlib. The
vendored definition of the entangled value has the strategy's two finite sets as
arbitrary types in the lowest universe, and the bridge reindexes them.

\effortHard\ \qlib

\begin{lemma}[Tensor powers of synchronous strategies]
\label{lem:tensor-power-pcc}
\uses{def:direct-repetition,def:sync-strategy,def:pcc,def:sync-value}
\lean{MIPRE.tensorFamily, MIPRE.SyncStrategy.tensorPow, MIPRE.SyncStrategy.value_tensorPow,
  MIPRE.SyncStrategy.isPCC_tensorPow, MIPRE.SyncStrategy.tensorPowDoubled,
  MIPRE.exists_perfectPCC_repeat_doubled}
\leanok
Let $\game$ be a synchronous game and $\strategy = (d, M)$ a synchronous strategy for it.
The $k$-fold tensor power $\strategy^{\ot k} = (d^k, M^{\ot k})$, with $M^{\bar{x}}_{\bar{a}}
= \bigotimes_{i=1}^{k} M^{x_i}_{a_i}$, is a synchronous strategy for the direct repetition
$\game^{k}$ whose value is the $k$-th power of the value of $\strategy$, and it is PCC if
$\strategy$ is. In particular a value-$1$ PCC strategy for the doubled game of a game
$\game$ gives one for the doubled game of $\game^{k}$, for every $k$, the tag being read on
every coordinate.
\end{lemma}
\begin{proof}
\leanok
The normalized trace on $\C^{d^k}$ of a tensor product is the product of the normalized
traces on $\C^d$ (the trace of a product of two projections being real), and the
distribution and predicate of $\game^{k}$ are products, so each summand of the value of
$\strategy^{\ot k}$ is the product over the coordinates of the corresponding summands of the
value of $\strategy$, and the sum over $k$-tuples of questions and answers is the $k$-th power
of the sum. Commutation is coordinatewise, on the support of the product distribution, which
is the product of the supports. For the doubled reading, the strategy on the doubled
repetition is the tensor power relabeled along $(t, \bar{x}) \mapsto ((t, x_i))_i$, which is
injective and covers the support of the doubled repeated distribution.
\end{proof}

\begin{lemma}[The repetition bound for answers of bounded length]
\label{lem:repetition-sound-bound}
\uses{thm:direct-repetition-q,def:normal-verifier}
\lean{MIPRE.Repetition.repConst, MIPRE.Repetition.card_answers_le,
  MIPRE.Repetition.quantumValue_repeat_le_soundBound}
\leanok
There is a universal constant $c' > 0$ such that for every game $\game$ whose answer
alphabets are the strings of length at most $B$, every $\eps > 0$ with $\valstar(\game) \leq
1 - \eps$ and every $k \geq 1$,
\begin{equation*}
  \valstar\bigl( \game^{k} \bigr) \;\leq\; \exp\Bigl( - \frac{c'\, \eps^{13}\, k}{B + 1} \Bigr)\;.
\end{equation*}
\end{lemma}
\begin{proof}
\leanok
Theorem~\ref{thm:direct-repetition-q} at $\eps' = 1 - \valstar(\game) \geq \eps$: there are at
most $3^{B+1}$ strings of length at most $B$, so $\log(\abs{\mA}\abs{\mB}) \leq 2(B + 1)
\log 3$, and $\eps' \leq 1 \leq B + 1$, so the denominator is at most $(B + 1)(1 + 2 \log 3)$;
the exponent is monotone in the gap, so the bound at $\eps'$ implies the one at $\eps$. Take
$c' = c / (1 + 2 \log 3)$. This is the form in which the soundness clause of
Theorem~\ref{thm:parallel-repetition} consumes the theorem.
\end{proof}

\begin{theorem}[Lin's tracial density theorem]
\label{thm:tracial-density}
\lean{MIPRE.Repetition.tracialDensity}
\leanok
\qlib{}
Let $\mX$ and $\mA$ be finite nonempty sets. Call a correlation $q(a, b \mid x, y)$ on
question alphabets $\mX, \mX$ and answer alphabets $\mA, \mA$ \emph{tracially
embeddable} if there are: a unital $*$-algebra $\alg$ over $\C$ with a tracial state
$\tau$ ($\tau(1) = 1$, $\tau(ab) = \tau(ba)$, $\tau(a^*) = \overline{\tau(a)}$); a
Hilbert space $\mH$ with a linear map $\iota : \alg \to \mH$ of dense range satisfying
$\langle \iota a, \iota b \rangle = \tau(a^* b)$, on which $\alg$ acts on the left by
bounded operators $L(a)$ with $L(a)\, \iota b = \iota(ab)$; a positive element $\sigma
\in \alg$ with $\tau(\sigma^* \sigma) = 1$; POVMs $\{E^x_a\}_{a} \subseteq \alg$ of
positive elements summing to $1$; and POVMs $\{G^y_b\}_{b}$ of positive bounded
operators on $\mH$ summing to $\Id$ and commuting with $L(\alg)$, such that
\begin{equation*}
  q(a, b \mid x, y) = \bigl\langle \iota\sigma,\ L(E^x_a)\, G^y_b\, \iota\sigma \bigr\rangle\;.
\end{equation*}
Then every commuting-operator correlation $p(a, b \mid x, y) = \bra{\psi} E^x_a F^y_b
\ket{\psi}$ on these alphabets (Definition~\ref{def:co-strategy-bipartite}) is an
$\ell^1$-limit of tracially embeddable correlations: for every $\delta > 0$ there is a
tracially embeddable $q$ with $\sum_{x, y, a, b} \abs{p(a, b \mid x, y) - q(a, b \mid
x, y)} < \delta$.
\end{theorem}
\begin{proof}
\leanok
This is Theorem~3.2 of~\cite{Lin23}, as formalized in~\cite{Vid26} (the density
programme of the vendored development, \texttt{Tracial/Density/Main.lean}). The
statement above is that of the vendored development, in its own vocabulary --- its
tracially embeddable correlations have the shape of Lin's Definition~3.1, with Bob's
measurement taken in the commutant of the left action; restating it through the
definitions of Section~\ref{ch:foundations} is pending.
\end{proof}

\paragraph{Source.} Direct parallel repetition for entangled strategies is Chapter~6 of
OpenAI's \emph{Ten advances}~\cite{OpenAI26}, with the Lean proof in the accompanying
repository; for commuting-operator strategies it is~\cite{Vid26}, whose proof reduces
to the tracially embeddable case by Theorem~\ref{thm:tracial-density} and then rounds
by adapting the argument of~\cite{OpenAI26}. Earlier value-form theorems cover
restricted classes of games or give polynomial rates; the anchored theorem
of~\cite{BVY17} is the one used in~\cite{JNVWY20}.
The rates $\eps^{13}$ and $\eps^{7}$ are those of the sources; any polynomial rate
suffices for the pipeline.

\paragraph{Comments.} Both vendored developments are sorry-free and use only the
standard axioms; the file \texttt{MIPRE/Background/Repetition/Axioms.lean} guards the
axioms of Theorem~\ref{thm:tracial-density} and of both bridged repetition theorems.
Vendoring
is done by \texttt{scripts/vendor-repetition.py}, never by hand; the copies are
read-only and the only local changes are recorded in each directory's
\texttt{README.md}. Nothing outside \texttt{MIPRE/Background/Repetition/} refers to
the vendored namespaces.

\begin{remark}[Direct versus anchored repetition]
\label{rem:direct-vs-anchored}
\uses{thm:direct-repetition-q,thm:direct-repetition-co}
Theorems~\ref{thm:direct-repetition-q} and~\ref{thm:direct-repetition-co} are
statements about values, while the anchored theorem of~\cite{BVY17} is a statement
about entanglement requirements: the anchored proof extracts from a good strategy for the repeated game a
strategy for the original game of \emph{no larger dimension}, whereas the proofs of the
direct theorems round through auxiliary resources (an embezzling state
in~\cite{OpenAI26}, a tracial approximation in~\cite{Vid26}) and give no control of the
dimension. Lin's proof of $\MIPco = \coRE$~\cite{Lin25} shows that the entanglement clauses
of~\cite{JNVWY20} are dispensable: with completeness and soundness of compression stated
in value form (perfect strategies are preserved, and $\valstar$ at most $\frac12$ is
preserved, a $\synval$ hypothesis side being converted by the free comparison
Lemma~\ref{lem:sync-le-valstar}), the halting reduction follows from the $\valstar$ half of
the enumeration of Lemma~\ref{lem:value-lower-approx}
and Kleene's recursion theorem alone (Lemma~\ref{lem:compressible-criterion}), and gap
amplification then needs only a value statement --- exactly the theorems of this section,
applied to the game as it is, without anchoring. The pipeline of
Section~\ref{ch:proof-structure} is stated in that form (Section~\ref{sec:departures}):
Theorem~\ref{thm:parallel-repetition} rests on Theorem~\ref{thm:direct-repetition-q},
and the anchored theorem with its toolkit --- the entropy inequalities, Uhlmann's
theorem, the quantum Raz lemma and Holenstein conditioning that its proof needs --- is
not on the path to the main theorem and is no longer stated in this blueprint; the
entanglement clauses are recorded in Remark~\ref{rem:entanglement-form}. The active
pipeline carries bipartite soundness throughout (Remark~\ref{rem:bipartite-route}); its
completeness witnesses use the separately specified synchronous PCC interface.
\end{remark}
